% Insert as A.43.  The construction in A.44 has the valuations in A.43.6.
\begin{statement}{A.43}{Proved}{The weight-three test on the twenty-four integral lifts}
Put $E_3=E_U/\Gamma_4\Pi$, where
$\Gamma_4\Pi=\overline{[\Pi,\Gamma_3\Pi]}$.  With the generators
and tangent of A.38, set
\[
 C=[x,y],\qquad D=[C,x],\qquad F=[C,y].
\]
The kernel of $E_3\to E_2$ is
$\widehat\Z(3)D\oplus\widehat\Z(3)F$.
For a class-two section $s(g)=(a(g),b(g),h(g);g)$ as in A.38,
write $\mathbf o_3(s)=(o_D(s),o_F(s))$ for its actual lifting
obstruction in $H^2_{\rm cts}(G_\Q,\widehat\Z(3))^2$.
Set $\varepsilon=(\chi-1)/2$ and $\eta=(\chi^2-1)/24$.
Its components are represented by the continuous two-cocycles
\[
 \begin{split}
 W_D&=-(a+\varepsilon)\cup h-\binom a2\cup b,\\
 W_F&=(b+\varepsilon)\cup(ab-h)+\binom b2\cup a-\eta\cup b.
 \end{split}
 \tag{A.43.1}
\]
Here $ab$ and $\binom a2$ are pointwise one-cochains with
coefficient weight two.  In particular all the expressions are
integral, including at two.  The twenty-four classes in A.40
have the exact obstruction values
\[
 \mathbf o_3(s+k\eta)=\mathbf o_3(s)
   -k\bigl((A+\kappa(-1))\cup[\eta],
           (B+\kappa(-1))\cup[\eta]\bigr),
                    \qquad k\in\Z/24.
 \tag{A.43.2}
\]

Let $(A,B)\in\mathcal G$, with ordinary coordinates $x_q$, and
let $\delta_q$ be its comparison class from (A.40.16).  For
$q\ne\ell$ denote the Frobenius coordinate of its $\ell$-part
by
\[
 d_{q,\ell}\in\Z_\ell/(q^2-1)\Z_\ell
       =H^1(G_{\Q_q},\Z_\ell(2)).
\]
Under the standard tame-residue isomorphism
\[
 \rho_{q,\ell}:H^2(G_{\Q_q},\Z_\ell(3))
       \xrightarrow{\ \sim\ }\Z_\ell/(q^2-1)\Z_\ell,
 \tag{A.43.3}
\]
the two localized weight-three obstructions are exactly
\[
 \rho_{q,\ell}\bigl(\mathbf o_3(s)|_{G_{\Q_q}}\bigr)
   =\bigl(-v_q(x_q)d_{q,\ell},
           -v_q(1-x_q)d_{q,\ell}\bigr).
 \tag{A.43.4}
\]
Thus their common zero locus at this place is the annihilator
of the two indicated valuation integers in the displayed finite
cyclic module.

Write $\mathscr L_i(A,B)$ for the geometric conjugacy classes
of $E_i$-sections with abelian image $(A,B)$, and $s_*$ for the
distinguished class of A.40.  On the ordinary locus with
\[
 (v_2(x_2),v_2(1-x_2))=(0,1),\qquad
 (v_3(x_3),v_3(1-x_3))=(1,0),
 \tag{A.43.5}
\]
the actual projection of section sets satisfies
\[
 \operatorname{im}\bigl(\mathscr L_3(A,B)\longrightarrow
                      \mathscr L_2(A,B)\bigr)\subseteq\{s_*\}.
 \tag{A.43.6}
\]
For the rational pair $(A,B)=(\kappa(3),\kappa(-2))$, this
image is exactly the class-two point-section class at $3$.
\end{statement}
\begin{proof}
\emph{The actual group, action and cochain representative.}
Commutator collecting generates $\Gamma_3\Pi/\Gamma_4\Pi$
by $D,F$.  To check their independence explicitly, send $x,y$
to $1+X,1+Y$ in the units of the noncommutative polynomial ring
truncated after degree three, with coefficients in
$\widehat\Z$.  This is a continuous homomorphism on the
class-three quotient.  The two commutators have expansions
\[
 \begin{split}
 D&\longmapsto 1+2XYX-YX^2-X^2Y,\\
 F&\longmapsto 1+XY^2-2YXY+Y^2X.
 \end{split}
 \tag{A.43.7}
\]
The coefficients of $X^2Y$ and $XY^2$ show that
$(u,v)\mapsto D^uF^v$ has zero kernel.  This proves the stated
basis of the kernel.  In particular every element of the geometric
class-three quotient has a unique form
\[
                         y^b x^a C^h D^uF^v.
 \tag{A.43.8}
\]
Its image in the Heisenberg model (A.38.2) is $(a,b,h)$;
the order of the first two factors is part of this identification.
Continuity follows from the profinite powers and the compact
coordinate spaces.

The arithmetic action from (A.38.6) has
$f_g\equiv C^{\eta(g)}\pmod{\Gamma_3\Pi}$, with this exact
coefficient and sign.  The coefficient calculation is
\cite[\S12.5.2, equation (12.35)]{BIR_Wickelgren2012}, using
the same tangent, generators and action
$g(y)=f_g^{-1}y^{\chi(g)}f_g$.
Writing $\lambda=\chi(g)$, the action on (A.43.8) is therefore
determined by
\[
 \begin{split}
 g(x)&=x^\lambda,\qquad
 g(y)=y^\lambda F^{-\lambda\eta(g)},\\
 g(C)&=C^{\lambda^2}
       D^{-\lambda\binom\lambda2}
       F^{-\lambda\binom\lambda2},\\
 g(D)&=D^{\lambda^3},\qquad g(F)=F^{\lambda^3}.
 \end{split}
 \tag{A.43.9}
\]
These formulas also verify the asserted cyclotomic action on
the kernel.  All binomial expressions are continuous
integer-valued functions on $\widehat\Z$.

The coordinate lift of $s$ to the geometric class-three quotient
is $\widetilde s(g)=y^{b(g)}x^{a(g)}C^{h(g)}$.  Its defect
\[
 \widetilde s(g)\,g(\widetilde s(t))\,
                       \widetilde s(gt)^{-1}
\]
has coordinates $(R_D,R_F)$ on the central kernel.  Applying
(A.43.9) and collecting the group words gives the exact cochain
identities
\[
 \begin{split}
 R_D&=W_D-\mathrm d(ah),\\
 R_F&=W_F+\mathrm d\bigl(a\binom b2-bh\bigr).
 \end{split}
 \tag{A.43.10}
\]
For completeness this collecting calculation can be performed
inside the same degree-three polynomial ring: use
\[
 \begin{split}
 (1+X)^a&=1+aX+\binom a2X^2+\binom a3X^3,\\
 C&=1+XY-YX+YX^2-XYX-XY^2+YXY,
 \end{split}
\]
and (A.43.7).  Substitute
$a(gt)=a(g)+\lambda a(t)$,
$b(gt)=b(g)+\lambda b(t)$ and
$h(gt)=h(g)+\lambda^2h(t)+a(g)\lambda b(t)$.
The degree-one and degree-two terms cancel.  The negative
$X^2Y$ coefficient and the $XY^2$ coefficient are the two
expressions in (A.43.10); these determine the central defect
by (A.43.7).  This proves (A.43.10) as an identity of
integer-valued polynomials, and continuity extends it to all
profinite exponents.  The symbolic program
\texttt{class3\_obstruction\_certificate.py} verifies precisely
these two identities over the rational polynomial ring, together
with (A.43.9).  Thus the computation does not discard a
finite-coefficient or continuous-cohomology correction.

The raw defect is a two-cocycle for the kernel action.  Formula
(A.43.10) proves that (A.43.1) consists of two-cocycles with the
same obstruction classes.  It is also the explicit comparison
with \cite[Proposition 17 and Theorem 19]{BIR_Wickelgren2012}:
the source's ordered coordinates $(b,a)_c=y^a x^b C^c$ become
our $(a,b,h)$ by $b_{\rm source}=a$,
$a_{\rm source}=b$, $c_{\rm source}=h$.
The obstruction is zero exactly on the image of the actual
map of section sets induced by $E_3\to E_2$: a continuous
primitive for the defect changes the last two coordinates of
the lift (A.43.8) to a homomorphism, and a section conversely
supplies such a primitive.

Replacing $h$ by $h+c$ for a continuous central cocycle gives,
by direct subtraction in (A.43.1),
\[
 \mathbf o_3(s+c)-\mathbf o_3(s)
       =-\bigl((A+\kappa(-1))\cup[c],
               (B+\kappa(-1))\cup[c]\bigr).
 \tag{A.43.11}
\]
Here $[\varepsilon]=\kappa(-1)$ by the compatible roots of
$-1$.  A.40 now proves (A.43.2), with no further choices of
defining systems for Massey products.

\emph{The exact local residue.}
The ordinary point section at $x_q$ has zero obstruction and
differs from $s|_{G_{\Q_q}}$ by $\delta_q$ after the unique
abelian alignment of A.40.  Applying (A.43.11) locally gives
\[
 \mathbf o_3(s)|_{G_{\Q_q}}
   =\bigl((A_q+\kappa_q(-1))\cup\delta_q,
           (B_q+\kappa_q(-1))\cup\delta_q\bigr).
 \tag{A.43.12}
\]
For $q\ne\ell$ every integral weight-two cocycle is unramified
by (A.40.6).  The tame Hochschild--Serre calculation identifies
$H^2(G_{\Q_q},\Z_\ell(3))$ with
$H^1(\operatorname{Gal}(\bar\F_q/\F_q),\Z_\ell(2))$:
prime-to-$\ell$ inertia has zero higher $\ell$-cohomology,
tame $\ell$-inertia has cohomological dimension one, and the
procyclic quotient has cohomological dimension one.  The same
calculation on finite coefficients, followed by inverse limit,
is valid for continuous integral cohomology since the finite
local $H^1$ groups are finite and their $\varprojlim^1$ is zero.
Frobenius evaluation now gives the isomorphism (A.43.3).

We fix its sign and normalization by the cochain residue from
A.27, now with coefficient weight three.  Write $T$ for a tame
inertia generator and $\Phi T\Phi^{-1}=T^q$.  For a normalized
two-cochain $v$, the residue is represented by
\[
 q^{-1}\left(v(\Phi,T)+v(\Phi T,\Phi^{-1})
       -v(\Phi,\Phi^{-1})
       -\sum_{i=1}^{q-1}v(T^i,T)\right).
 \tag{A.43.13}
\]
The telescoping computation of A.27 gives
$(q^2-1)u(T)$ on $v=\mathrm du$; this is the differential
appropriate to weight three.  The discarded inertia is handled
by exactly the coefficient-compatible averaging map of A.27.

Let $w$ be a weight-one cocycle, with $w(T)=t$ and
$w(\Phi)=\alpha$, and let $c$ be an unramified weight-two
cocycle with $c(\Phi)=d$.  In (A.43.13), the cochain
$v=w\cup c$ has first and last terms zero and middle
difference
\[
                   -\bigl(\alpha+qt\bigr)d+\alpha d=-qtd.
\]
Its residue is therefore $-td$ modulo $(q^2-1)\Z_\ell$.
The inertia values of the two weight-one Kummer classes in
(A.43.12) are precisely $v_q(x_q)$ and $v_q(1-x_q)$;
$\kappa_q(-1)$ has inertia value zero.  This proves (A.43.4)
with the exact comparison map and sign.

\emph{The two fixed valuations.}
For a class in the image of $\mathscr L_3(A,B)$, both
obstructions vanish.  Under (A.43.5), (A.43.4) at
$(q,\ell)=(2,3)$ forces $d_{2,3}=0$, and at $(3,2)$ it
forces $d_{3,2}=0$.  The isomorphism (A.40.4) says that
exactly one global class-two lift has these two comparisons
zero, namely $s_*$.  This proves (A.43.6).

For $x=3$, the full rational point section exists and gives
one element in this image.  More explicitly, take its class-two
class as $s_0$.  For $s_0+k\eta$, the two local obstruction
vectors at the indicated coefficient pairs are
\[
       (0,k/8)\in(\Z/3)^2,
       \qquad(k/3,0)\in(\Z/8)^2.
 \tag{A.43.14}
\]
They vanish together exactly for $k=0$ in $\Z/24$.
This identifies the image as the point-section class itself.
\end{proof}

The reduction of the integral central adjustment map at two
is especially small:
\[
 \operatorname{im}\bigl(H^1(G_\Q,\widehat\Z(2))
                \longrightarrow H^1(G_\Q,\Z/2)\bigr)
                      =\{0,\kappa_2(2)\}.
 \tag{A.43.15}
\]
Indeed $\eta\bmod2$ is one exactly when
$\chi_2\equiv\pm3\pmod8$, and zero when
$\chi_2\equiv\pm1\pmod8$.  This is the quadratic character of
$\Q(\sqrt2)$, as also computed in
\cite[Lemma 31]{BIR_Wickelgren2012}.  Formula (A.43.15)
specifies exactly which mod-two changes of central defining
system come from the continuous integral lift fiber; the
finite-coefficient choice is required to lie in this image.

The nonrational pairs constructed in A.44 have the valuations
(A.43.5).  Their possible class-three projections are therefore
confined to the single class $s_*$.  No vanishing of
$\mathbf o_3(s_*)$ for those pairs is proved here.  Equations
(A.43.1), (A.43.4) and the actual extension $E_3\to E_2$
specify the remaining calculation, including its continuous
global cohomology class.
